% =====================================================================
%  SIGLAS Y ABREVIATURAS — paquete acronym
%
%  Defina cada sigla aquí, en orden alfabético, y úsela en el texto con
%    \ac{INEGI}   -> la 1.ª vez: "Instituto Nacional de ... (INEGI)";
%                    después solo "INEGI"
%    \acs{INEGI}  -> siempre la sigla      \acl{INEGI} -> siempre el nombre
%  Si quita esta lista de main.tex, reemplace también los \ac del texto.
% =====================================================================
\chapter*{Siglas y abreviaturas}
\addcontentsline{toc}{chapter}{Siglas y abreviaturas}
\markboth{Siglas y abreviaturas}{Siglas y abreviaturas}

% El argumento opcional es la sigla más larga (alinea las descripciones).
\begin{acronym}[CONEVAL]
  \acro{CONEVAL}{Consejo Nacional de Evaluación de la Política de Desarrollo Social}
  \acro{INEGI}{Instituto Nacional de Estadística y Geografía}
  \acro{MCO}{mínimos cuadrados ordinarios}
  \acro{PIB}{producto interno bruto}
\end{acronym}
